\paragraph{Classical argument supervision.}
We used release 2.16 of the Kyoto Classical Chinese Universal Dependencies training corpus. We extracted verbal subject and object relations whose argument head was labelled as a personal proper noun. Contiguous flat name components were retained. We supplied a bilingual question wording for each eligible predicate and grammatical role. The training set contained 194 original clauses with two annotated arguments of one verb, and 1,979 constructed pages joining clauses with different verbs and distinct people in the same grammatical role. The 3,960 eligible source facts covered 49 verbs used in training. Counterfactual pages exchanged the two complete name strings; the order of joined clauses was varied across examples. The newly reserved 16 excerpts from the \textit{Zuo Tradition} have independently written questions in both languages. Neither classical evaluation supplied names or candidates to the reader. Excerpts were excluded from training after Unicode normalization and removal of punctuation. The smaller models shared their 64,666 parameters and initial weights with the previous distilled reader. The aligned reader minimized the mean of the full-page and two-name negative log likelihoods of the annotated response image; all occurrences contributed. The independent-label fit averaged the same loss over both name images for every question, keeping the page, labels and optimizer schedule fixed. The registered resource limit was one GPU and 4 GiB allocated for each fit. Position alone can follow a name exchange if the question already identifies the required position, as demonstrated in Lean. These contrasts therefore measure a specific behaviour in the context of the actual question and controls.
